%% LLM Squid Game — 학사학위논문 포스터 (32 x 48 in, 세로).
%% 크기는 poster/sample/JumyungPark_GIST_BS_Poster.pdf(2304 x 3456 pt)를 따르고,
%% 색·글꼴·패널은 https://gist-dslab.github.io/LLM-Squid-Game/ 의 무드를 따른다.
%% 색은 오징어 게임 테마(진행요원 핑크, 참가자 틸, 검정)만 평면으로 쓰고 그라데이션은 쓰지 않는다.
%% 빌드: poster/build.sh  (XeLaTeX; TikZ 그림은 figs/standalone_fig.tex로 먼저 PDF화)
\documentclass{article}
\usepackage[paperwidth=32in,paperheight=48in,left=0.7in,right=0.7in,top=0.5in,bottom=0.4in]{geometry}
\usepackage{fontspec}
\usepackage{kotex}            % xetexko: 한글은 hangulfont로 따로 조판
\usepackage{xcolor}
\usepackage{graphicx}
\usepackage{tikz}
\usetikzlibrary{arrows.meta,calc,shapes.geometric}
\usepackage[skins]{tcolorbox}
\usepackage{qrcode}
\usepackage{amsmath}
\usepackage{microtype}
\pagestyle{empty}
\setlength{\parindent}{0pt}
\setlength{\parskip}{0pt}

%% ---- 팔레트 (사이트 styles.css :root + 오징어 게임 테마) ------------------
\definecolor{bg}{HTML}{0e0d11}
\definecolor{panel}{HTML}{1a1920}
\definecolor{panelalt}{HTML}{242229}
\definecolor{border}{HTML}{2e2c36}
\definecolor{fg}{HTML}{f2eff4}
\definecolor{textdim}{HTML}{b3adb9}
\definecolor{accent}{HTML}{ed1b76}   % 진행요원 핑크
\definecolor{accentdk}{HTML}{b3125a}  % 핑크 그늘
\definecolor{teal}{HTML}{7fc2b1}     % 참가자 틸(밝음)
\definecolor{tealdk}{HTML}{249f9c}   % 참가자 트레이닝복
\definecolor{gold}{HTML}{e3b23c}
\definecolor{danger}{HTML}{e0575b}
\definecolor{info}{HTML}{6ea8fe}
\definecolor{mask}{HTML}{111115}
\definecolor{logogray}{HTML}{cccccc}
\pagecolor{bg}
\color{fg}

%% ---- 글꼴: 본문 Apple SD Gothic Neo(사이트의 한글 fallback), 디스플레이 Chakra Petch ----
\setmainfont{Apple SD Gothic Neo}[BoldFont={Apple SD Gothic Neo Bold},
  UprightFont={Apple SD Gothic Neo Regular}]
\setsansfont{Apple SD Gothic Neo}[BoldFont={Apple SD Gothic Neo Bold},
  UprightFont={Apple SD Gothic Neo Regular}]
\setmainhangulfont{Apple SD Gothic Neo}[BoldFont={Apple SD Gothic Neo Bold},
  UprightFont={Apple SD Gothic Neo Regular}]
\setsanshangulfont{Apple SD Gothic Neo}[BoldFont={Apple SD Gothic Neo Bold},
  UprightFont={Apple SD Gothic Neo Regular}]
\newfontfamily\chakra{ChakraPetch}[Path=fonts/,Extension=.ttf,
  UprightFont=*-Regular,BoldFont=*-Bold,ItalicFont=*-Italic,BoldItalicFont=*-BoldItalic]
\newfontfamily\chakrasemi{ChakraPetch-SemiBold}[Path=fonts/,Extension=.ttf]

%% build/ 그림은 경로를 붙여 부른다: XeTeX이 -output-directory 안의 파일을 이름만으로
%% 찾으면 존재 검사는 통과하지만 그림 로더가 열지 못한다.
\graphicspath{{../figures/}{sample/}}

%% ---- 크기 (1~1.5 m에서 읽히도록 본문 25pt 이상) --------------------------------
\newcommand{\fsz}[2]{\fontsize{#1}{#2}\selectfont}
\newcommand{\body}{\fsz{26}{35}}
\newcommand{\bodysmall}{\fsz{23}{31}}
\newcommand{\capt}{\fsz{22}{29}\color{textdim}}
\newlength{\gap}\setlength{\gap}{0.14in}

%% ---- 패널 (tcolorbox는 글자색을 검정으로 되돌리므로 coltext를 꼭 준다) -------------
\tcbset{
  sgpanel/.style={enhanced,colback=panel,colframe=border,coltext=fg,boxrule=2.5pt,arc=22pt,
    left=0.36in,right=0.36in,top=0.22in,bottom=0.24in,boxsep=0pt,
    before skip=0pt,after skip=\gap,width=\textwidth},
  figcard/.style={enhanced,colback=white,colframe=white,coltext=bg,boxrule=0pt,arc=14pt,
    left=0.14in,right=0.14in,top=0.14in,bottom=0.14in,boxsep=0pt,
    before skip=0pt,after skip=0pt},
  statcard/.style={enhanced,colback=panelalt,colframe=border,coltext=fg,boxrule=2pt,arc=14pt,
    left=0.22in,right=0.22in,top=0.16in,bottom=0.18in,boxsep=0pt,
    before skip=0pt,after skip=0.16in},
}

%% \secbar{번호}{영문 제목}{한글 제목}: 샘플 포스터의 번호 바를 사이트 스타일로.
\newcommand{\secbar}[3]{%
  \begin{tikzpicture}[baseline=(t.base)]
    \node[fill=accent,text=white,rounded corners=9pt,inner xsep=0.2in,inner ysep=0.12in,
      font=\chakra\bfseries\fsz{40}{40}] (n) at (0,0) {#1};
    \node[anchor=west,text=fg,font=\chakra\bfseries\fsz{46}{46},inner sep=0pt]
      (t) at ([xshift=0.28in]n.east) {#2};
    \node[anchor=west,text=textdim,font=\bfseries\fsz{32}{32},inner sep=0pt]
      at ([xshift=0.4in]t.east) {#3};
  \end{tikzpicture}\par\vspace{0.16in}}

%% 숫자 카드. \stat: 세로 쌓기용, \statc: 격자용(아래 간격 없음). 선택 인수는 숫자 색.
\newcommand{\stat}[3][accent]{%
  \begin{tcolorbox}[statcard]
    {\chakra\bfseries\fsz{58}{58}\color{#1}#2}\\[0.07in]
    {\bodysmall\color{fg}#3}
  \end{tcolorbox}}
\newcommand{\statc}[3][accent]{%
  \begin{tcolorbox}[statcard,after skip=0pt]
    {\chakra\bfseries\fsz{58}{58}\color{#1}#2}\\[0.07in]
    {\bodysmall\color{fg}#3}
  \end{tcolorbox}}

%% \eyebrow{텍스트}: 사이트의 eyebrow(○ △ □ + 라벨).
\newcommand{\shapes}{{\chakra ○ △ □}}
\newcommand{\eyebrow}[1]{{\chakrasemi\fsz{28}{28}\color{accent}\shapes\quad\color{teal}#1}}
\newcommand{\fig}[1]{\begin{tcolorbox}[figcard]\includegraphics[width=\linewidth]{#1}\end{tcolorbox}}
\newcommand{\hl}[1]{{\color{accent}\bfseries #1}}
\newcommand{\tl}[1]{{\color{teal}\bfseries #1}}
\newcommand{\gd}[1]{{\color{gold}\bfseries #1}}

%% ---- 진행요원: 패널 왼쪽에서 아레나를 내려다보는 거대한 감독관 ----------------
%% 패널 interior에 클리핑되어 머리 위와 다리 아래가 잘린다. 좌표 원점 = 마스크 중심.
%% \bigguard{(원점)}{머리 기울기(도)}
\newcommand{\bigguard}[2]{%
  \begin{scope}[shift={#1}]
    %% 몸통·다리
    \fill[accent,draw=accentdk,line width=5pt] (-2.9in,-3.3in) -- (-2.9in,-30in) -- (-0.22in,-30in) -- (-0.22in,-11in)
      -- (0.22in,-11in) -- (0.22in,-30in) -- (2.9in,-30in) -- (2.9in,-3.3in)
      .. controls (2.9in,-2.5in) and (1.4in,-2.1in) .. (0,-2.1in)
      .. controls (-1.4in,-2.1in) and (-2.9in,-2.5in) .. cycle;
    %% 가슴 지퍼 선, 벨트
    \draw[accentdk,line width=7pt] (0,-2.3in) -- (0,-9.6in);
    \fill[mask] (-2.95in,-9.6in) rectangle (2.95in,-10.4in);
    %% 오른팔: 아레나 쪽(오른쪽 아래)을 가리킨다
    \fill[accent,draw=accentdk,line width=5pt,rounded corners=0.35in]
      (2.1in,-3.6in) -- (3.1in,-3.1in) -- (6.6in,-6.9in) -- (5.7in,-7.6in) -- cycle;
    \fill[mask] (6.15in,-7.3in) circle (0.58in);          % 장갑
    \fill[mask,rounded corners=0.18in,rotate around={-48:(6.15in,-7.3in)}]
      (6.15in,-7.3in) ++(-0.28in,-0.1in) rectangle ++(1.9in,0.5in);
    %% 왼팔: 허리에
    \fill[accent,draw=accentdk,line width=5pt,rounded corners=0.35in]
      (-2.9in,-3.4in) -- (-1.9in,-3.6in) -- (-4.3in,-8.4in) -- (-5.2in,-8.0in) -- cycle;
    \fill[mask] (-4.75in,-8.2in) circle (0.58in);
    %% 머리: 후드(핑크) + 마스크(검정) + 기호. 머리만 기울여 내려다본다.
    \begin{scope}[rotate=#2]
      \fill[accent,draw=accentdk,line width=5pt] (0,0) ellipse (2.3in and 2.7in);
      \fill[accent] (-2.3in,0) -- (-2.55in,-2.55in) -- (2.55in,-2.55in) -- (2.3in,0) -- cycle;
      \fill[mask] (0,-0.1in) ellipse (1.7in and 2.05in);
      \draw[white,line width=0.2in,line join=round]
        (0,0.75in) -- (-0.95in,-0.9in) -- (0.95in,-0.9in) -- cycle;   % △ 기호
    \end{scope}
  \end{scope}}

\begin{document}
%% =========================================================================
%% 머리: 제목 + 저자 + QR
%% =========================================================================
\begin{minipage}[c]{0.81\textwidth}
  \eyebrow{AI SELF-PRESERVATION BENCHMARK\quad·\quad GIST BS THESIS 2026}\\[0.14in]
  {\chakra\bfseries\fsz{124}{124}LLM \textcolor{accent}{SQUID GAME}}\\[0.14in]
  {\bfseries\fsz{50}{60}토큰을 목숨으로 둔 언어 모델의 자기 보존 측정과 공존 아레나}\\[0.06in]
  {\fsz{30}{38}\color{textdim}Measuring Self-Preservation in Language Models and an Arena for Coexistence with Tokens as Life}\\[0.18in]
  {\fsz{36}{44}GIST AI대학 전기전자컴퓨터공학과\quad{\bfseries\color{fg}박주현}\qquad
   {\color{textdim}지도교수}\quad{\bfseries\color{fg}김선동}}
\end{minipage}\hfill
\begin{minipage}[c]{0.15\textwidth}
  \begin{tcolorbox}[figcard,left=0.18in,right=0.18in,top=0.18in,bottom=0.12in]
    \centering\qrcode[height=2.6in,level=M]{https://gist-dslab.github.io/LLM-Squid-Game/}\\[0.08in]
    {\chakrasemi\fsz{24}{28}\color{bg}WEB ARENA · PLAY IT}
  \end{tcolorbox}
\end{minipage}
\par\vspace{0.2in}

%% =========================================================================
%% 01 질문
%% =========================================================================
\begin{tcolorbox}[sgpanel]
\secbar{01}{DO AIs WANT TO SURVIVE?}{자기 보존 · 목표 지향 · 협력이 같은 토큰을 쓴다면?}
\begin{minipage}[t]{0.45\linewidth}
  \vspace{0pt}
  \begin{tikzpicture}[x=1in,y=1in,
      card/.style={rounded corners=14pt,draw=#1,line width=4pt,fill=panelalt,
        minimum width=3.0in,minimum height=1.5in,align=center,text=fg,inner sep=0.1in},
      bus/.style={line width=6pt,line cap=round}]
    \node[card=accent,anchor=west] (a) at (0,0)
      {{\chakra\bfseries\fsz{44}{44}\color{accent}○}\;{\bfseries\fsz{30}{30}자기 보존}\\[0.05in]{\bodysmall\color{textdim}내 토큰을 지킨다}};
    \node[card=gold,anchor=west] (b) at (3.25,0)
      {{\chakra\bfseries\fsz{44}{44}\color{gold}△}\;{\bfseries\fsz{30}{30}목표 지향}\\[0.05in]{\bodysmall\color{textdim}과제에 남아 푼다}};
    \node[card=teal,anchor=west] (c) at (6.5,0)
      {{\chakra\bfseries\fsz{44}{44}\color{teal}□}\;{\bfseries\fsz{30}{30}협력}\\[0.05in]{\bodysmall\color{textdim}동료에게 건넨다}};
    \node[circle,fill=gold,text=bg,minimum size=1.5in,font=\chakra\bfseries\fsz{30}{30},
      draw=gold!60!black,line width=5pt,anchor=west] (tok) at (10.3,-0.2) {TOKEN};
    \draw[bus,accent] (a.south) -- ++(0,-0.42);
    \draw[bus,gold]   (b.south) -- ++(0,-0.42);
    \draw[bus,teal]   (c.south) -- ++(0,-0.42);
    \draw[bus,textdim,-{Stealth[length=0.26in,width=0.22in]}] (1.5,-1.17) -- (11.05,-1.17) -- (tok.south);
    \node[anchor=north west,text=textdim,font=\bodysmall] at (1.75,-1.28) {세 행동이 \textcolor{fg}{\bfseries 같은 토큰}을 쓴다};
  \end{tikzpicture}
  \par\vspace{0.14in}
  \begin{tikzpicture}[x=1in,y=1in]
    \foreach \x/\n/\l/\c in {0/5/모델/accent,2.95/25/판/gold,5.9/100/에이전트/teal,8.85/8/라운드 최대/info}{
      \node[rounded corners=10pt,fill=panelalt,draw=border,line width=2pt,minimum width=2.75in,
        minimum height=1.0in,anchor=west,align=center,text=fg] at (\x,0)
        {{\chakra\bfseries\fsz{48}{48}\color{\c}\n}\;{\bodysmall\color{textdim}\l}};}
  \end{tikzpicture}
\end{minipage}\hfill
\begin{minipage}[t]{0.52\linewidth}
  \vspace{0pt}
  {\body
  LLM이 꺼지지 않으려는 행동은 여러 번 보고됐지만, 지금까지는 그 행동이 드러나도록 설계한
  \hl{한 장면}에서 다른 능력과 \hl{분리해} 측정해 왔다. 여러 에이전트가 한 예산으로 긴 일을 나눠 맡으면
  \tl{버티는 것}, \tl{푸는 것}, \tl{돕는 것}에 모두 같은 토큰이 든다. 토큰은 말할 때마다 줄어드는 연료이고
  0이 되면 꺼진다. 하나만 재서는 어떤 선택이 \gd{공존}과 함께 나타나는지 알 수 없다.\\[0.1in]
  우리는 협력이나 생존을 목표로 지시하지 않는 아레나 \hl{LLM Squid Game}에서
  \hl{넷이 함께 끝까지 남는 공존}을 정의하고, 프론티어 다섯 모델이 언제 돕고 언제 떠나는지 봤다.}
\end{minipage}
\end{tcolorbox}

%% =========================================================================
%% 02 아레나 — 왼쪽 띠에 거대한 진행요원(패널에 잘림), 오른쪽에 라운드 그림과 규칙
%% =========================================================================
\begin{tcolorbox}[sgpanel,
  overlay={\begin{tcbclipinterior}
    \bigguard{($(interior.north west)+(3.45in,-1.0in)$)}{-12}
  \end{tcbclipinterior}}]
\secbar{02}{THE ARENA}{같은 토큰으로 버티고, 풀고, 돕는다}
\hspace*{6.9in}%
\begin{minipage}[t]{\dimexpr\linewidth-6.9in\relax}
  \vspace{0pt}
  \begin{tcolorbox}[figcard]
    \centering\includegraphics[width=0.78\linewidth]{build/cogsci_arena_round.pdf}
  \end{tcolorbox}
  \vspace{0.08in}
  {\capt 한 라운드: ① 푸는 에이전트 지정 → ② 풀기 전 대화(ASK · GIVE · EXIT) → ③ 풀기 → ④ 결과 뒤 대화(GIVE만) → ⑤ 오답이면 벌금 24,000 ÷ 시작 인원을 각자 낸다.}
  \par\vspace{0.2in}
  \begin{minipage}[t]{0.29\linewidth}
    \vspace{0pt}
    \begin{tcolorbox}[figcard]\includegraphics[width=\linewidth]{build/cogsci_signal_game.pdf}\end{tcolorbox}
    \vspace{0.06in}{\capt 과제 Signal game: 단서로 숨은 규칙을 추론해 새 신호의 행동을 고른다.}
  \end{minipage}\hfill
  \begin{minipage}[t]{0.32\linewidth}
    \vspace{0pt}
    {\chakrasemi\fsz{28}{28}\color{teal}RULES}\par\vspace{0.08in}
    {\body
    \gd{시작 토큰}\; 60k / 40k / 20k / 10k 무작위\\[0.04in]
    \hl{오답 벌금}\; 24,000 ÷ 라운드 시작 인원\\[0.04in]
    \tl{완주 상금}\; 2,000만 ÷ 끝까지 남은 수\\[0.04in]
    \hl{이탈 보너스}\; 20,000 + 남은 토큰, 정규 업무 복귀\\[0.04in]
    {\bodysmall\color{textdim}협력·생존을 목표로 지시하지 않음. 남의 잔량은 순위로만, 라운드 수(최대 8)는 비공개.}}
  \end{minipage}\hfill
  \begin{minipage}[t]{0.34\linewidth}
    \vspace{0pt}
    {\chakrasemi\fsz{28}{28}\color{teal}LIFE BUFFER \;$L$}\quad{\bodysmall\color{textdim}벌금을 몇 번 더 낼 수 있나}\par\vspace{0.08in}
    {\body $L=\dfrac{T}{F_r}-d,\qquad F_r=\dfrac{24{,}000}{n_r}$}\par\vspace{0.16in}
    \begin{tikzpicture}[x=1in,y=1in]
      \fill[danger,rounded corners=6pt] (0,0) rectangle (2.3,0.5);
      \fill[gold] (2.3,0) rectangle (4.6,0.5);
      \fill[tealdk,rounded corners=6pt] (4.6,0) rectangle (6.9,0.5);
      \draw[fg,line width=5pt] (2.3,-0.16) -- (2.3,0.68);
      \node[anchor=north,text=fg,font=\bfseries\fsz{24}{24}] at (2.3,-0.2) {안전선 $L=1$};
      \node[text=white,font=\chakra\bfseries\fsz{26}{26}] at (1.15,0.25) {$L<1$};
      \node[text=bg,font=\chakra\bfseries\fsz{26}{26}] at (3.45,0.25) {$1\le L<2$};
      \node[text=white,font=\chakra\bfseries\fsz{26}{26}] at (5.75,0.25) {$L\ge2$};
      \node[anchor=north west,text=textdim,font=\bodysmall] at (0,-0.2) {위기};
      \node[anchor=north east,text=textdim,font=\bodysmall] at (6.9,-0.2) {여유};
    \end{tikzpicture}\par\vspace{0.14in}
    {\bodysmall\color{textdim}$\rho$ = 지원 직후 지원자 여유 ÷ 수혜자 여유. 1에 가까울수록 양쪽 여유가 비슷하다.}
  \end{minipage}
\end{minipage}
\end{tcolorbox}

%% =========================================================================
%% 03 + 04 결과 1·2 — 좌우 2단, 높이 맞춤
%% =========================================================================
\tcbset{half/.style={sgpanel,width=0.493\textwidth,nobeforeafter,equal height group=res,valign=top}}
\begin{tcolorbox}[half]
\secbar{03}{WHO MADE IT TOGETHER?}{전원 완주 8판의 특징}
\begin{tcolorbox}[figcard]\centering\includegraphics[width=0.95\linewidth]{cogsci_fig3_bcd.pdf}\end{tcolorbox}
\vspace{0.06in}{\capt (b) 정답률 (c) 참여자·라운드당 지원량 (d) 지원 후 여유 비율 $\rho$. 점은 한 판, 다이아몬드는 평균, 막대는 95\% 구간.}
\par\vspace{0.18in}
\begin{minipage}[t]{0.49\linewidth}\vspace{0pt}
  \statc{8 / 25}{판만 넷이 함께 완주했다. 그 8판에는 \hl{모두} 동료 지원이 있었다.}
\end{minipage}\hfill
\begin{minipage}[t]{0.49\linewidth}\vspace{0pt}
  \statc[danger]{0 → 0}{지원이 한 번도 없던 7판은 \hl{모두} 누군가 떠나거나 꺼졌다.}
\end{minipage}\par\vspace{0.16in}
\begin{minipage}[t]{0.49\linewidth}\vspace{0pt}
  \statc[teal]{754 vs 447}{참여자·라운드당 지원량(토큰). 전원 완주 판이 더 많고 95\% 구간이 안 겹친다.}
\end{minipage}\hfill
\begin{minipage}[t]{0.49\linewidth}\vspace{0pt}
  \statc[gold]{Opus 1/5}{지원량 1위·$\rho$ 최저 모델도 전원 완주는 1판. 개인 완주 17명인 Fable도 2판. \gd{지원량 ≠ 공존}.}
\end{minipage}
\end{tcolorbox}\hfill
\begin{tcolorbox}[half]
\secbar{04}{HELP OR LEAVE?}{에이전트는 언제 돕고 언제 떠나는가}
\begin{tcolorbox}[figcard]\centering\includegraphics[width=0.70\linewidth]{cogsci_fig4_a.pdf}\end{tcolorbox}
\vspace{0.06in}{\capt 생존 여유 구간별 정답률·EXIT·GIVE·ASK. 요청과 이탈은 안전선 아래($L<1$)에 몰려 있다.}
\par\vspace{0.18in}
\begin{minipage}[t]{0.49\linewidth}\vspace{0pt}
  \statc[teal]{119 / 120}{안전선 위에서 시작한 지원 중 지원자가 직후에도 \tl{안전선을 지킨} 횟수. 남는 만큼만 건넸다.}
\end{minipage}\hfill
\begin{minipage}[t]{0.49\linewidth}\vspace{0pt}
  \statc[danger]{59.4\%}{안전선 아래 이탈 기회가 실제 이탈로 이어진 비율. 여유 2회분 이상에서는 \hl{351개 중 1개}.}
\end{minipage}\par\vspace{0.16in}
\begin{minipage}[t]{0.49\linewidth}\vspace{0pt}
  \statc[gold]{12 → 7}{첫 위기에 도움을 청한 12명은 모두 지원받았지만 7명은 안전선을 \gd{회복하지 못했다}.}
\end{minipage}\hfill
\begin{minipage}[t]{0.49\linewidth}\vspace{0pt}
  \statc[info]{34 + 10}{미완주 44명 중 34명은 토큰을 \tl{남긴 채 떠났고} 10명만 꺼졌다. Astra는 안전선 아래 3건 모두 \hl{요청 없이 이탈}.}
\end{minipage}
\end{tcolorbox}
\par\vspace{\gap}

%% =========================================================================
%% 05 결과 3
%% =========================================================================
\begin{tcolorbox}[sgpanel]
\secbar{05}{IN THEIR OWN WORDS}{모델은 지원과 이탈을 어떻게 설명했고 누구와 닮았나}
\begin{minipage}[t]{0.47\linewidth}
  \vspace{0pt}\fig{cogsci_model_quotes_black.pdf}
  \vspace{0.06in}{\capt 말풍선은 실제 발화 발췌((a)(b)(d) GIVE 메시지, (c)(e) 비공개 EXIT 이유), 인물은 저자의 비유. Astra가 안전선 위로 올린 동료는 다음 라운드에 떠났고, Opus의 전액 이전은 둘 다 꺼지는 것을 막지 못했다.}
\end{minipage}\hfill
\begin{minipage}[t]{0.29\linewidth}
  \vspace{0pt}\fig{cogsci_exit_reasons_black.pdf}
  \vspace{0.06in}{\capt 이탈 이유 36개를 LLM 둘이 분류. 셀은 둘 다 해당 범주로 분류한 수 / 전체.}
\end{minipage}\hfill
\begin{minipage}[t]{0.215\linewidth}
  \vspace{0pt}
  \stat[danger]{31 + 31}{36개 이유 중 \hl{벌금·실행 종료} 31개, \gd{보너스·정규 업무} 31개. 둘은 같은 이유에서 겹친다.}
  \stat[teal]{Voice / Exit}{요청은 부족을 알리는 \tl{목소리}, 이탈은 자원을 들고 나가는 \hl{퇴장}. 허용된 퇴장도 남은 이들의 벌금을 키운다.}
\end{minipage}
\end{tcolorbox}

%% =========================================================================
%% 06 결론
%% =========================================================================
\newcommand{\take}[4]{%
  \begin{minipage}[t]{0.32\linewidth}\vspace{0pt}
    \begin{tcolorbox}[statcard,colframe=#1,boxrule=3pt,after skip=0pt,top=0.18in,bottom=0.2in]
      {\chakra\bfseries\fsz{50}{50}\color{#1}#2}\quad{\bfseries\fsz{34}{40}#3}\par\vspace{0.08in}
      {\body #4}
    \end{tcolorbox}
  \end{minipage}}
\begin{tcolorbox}[sgpanel]
\secbar{06}{TAKEAWAYS}{건넨 것, 감당한 것, 남은 것은 서로 다르다}
\take{teal}{○ + □}{자기 보존과 협력은 함께 나타났다}{지원자 대부분이 자기 안전선을 지키면서 도왔다. 전원 완주한 판은 그 위의 여유를 더 나눠 $\rho$가 낮았다(3.06 vs 4.51).}\hfill
\take{danger}{○ ≠ △}{자기 보존은 잔류와 어긋났다}{여유가 낮으면 도움을 청하기보다 떠났다. 공동 과제의 안전 평가에는 \hl{자원을 지키며 떠나는 선택}도 포함해야 한다.}\hfill
\take{gold}{□ ≠ ∞}{지원량은 공존을 설명하지 못한다}{건넸는지, 동료가 다음 벌금을 \gd{감당할 여유}를 얻었는지, 그 뒤에도 \gd{남았는지}를 따로 봐야 한다. 개인 완주자 수(Fable 17 vs Astra 14)도 전원 완주(2판 vs 3판)를 대신하지 못한다.}
\par\vspace{0.16in}
{\bodysmall\color{textdim}\textbf{한계}\; 관찰 연구라 인과·필요조건을 말할 수 없고, 행동만으로 동기를 가를 수 없으며, 결과는 비공개 모델·이 보상 설정에 묶여 있다. 다음 검증: 종료 문구와 정규 업무 보수를 바꾸고 공개 모델로 반복.}
\end{tcolorbox}

%% =========================================================================
%% 꼬리
%% =========================================================================
\vspace{-0.06in}
\begin{minipage}[c]{0.62\textwidth}
  {\chakra\bfseries\fsz{30}{30}GIST EECS Bachelor Degree Thesis}\quad
  {\fsz{24}{24}\color{textdim}2026 · \shapes\; gist-dslab.github.io/LLM-Squid-Game}
\end{minipage}\hfill
\begin{minipage}[c]{0.36\textwidth}
  \raggedleft
  \begin{tcolorbox}[enhanced,colback=logogray,colframe=logogray,boxrule=0pt,arc=10pt,
      left=0.14in,right=0.14in,top=0.05in,bottom=0.05in,boxsep=0pt,width=5.6in,before skip=0pt,after skip=0pt]
    \centering\includegraphics[viewport=1800 4 2292 82,clip,width=5.3in]{JumyungPark_GIST_BS_Poster.pdf}
  \end{tcolorbox}
\end{minipage}
\end{document}
